\documentclass[11pt]{article}
\usepackage[T1]{fontenc}
\usepackage{lmodern}
\usepackage{amsmath,amssymb,amsthm}
\usepackage[margin=1in]{geometry}
\usepackage[colorlinks=true,linkcolor=blue,citecolor=blue,urlcolor=blue]{hyperref}

\newtheorem{theorem}{Theorem}
\newtheorem{corollary}{Corollary}
\newtheorem{lemma}{Lemma}

\newcommand{\C}{\mathbb C}
\newcommand{\N}{\mathbb N}
\newcommand{\ellp}{\ell^p}
\newcommand{\supp}{\operatorname{supp}}
\newcommand{\spn}{\operatorname{span}}

\title{Partial Result: Amenable-Radical Obstruction for Reduced \texorpdfstring{$L^p$}{Lp} Group Banach Algebras}
\author{Run \texttt{fa\_banach\_001}, agent \texttt{agent\_super\_01}}
\date{2026-07-01}

\begin{document}
\maketitle

\section*{Status}

This is a partial result for the reverse-implication question in Phillips,
\emph{Simplicity of reduced group Banach algebras}, arXiv:1909.11278. Phillips
asks whether simplicity of $F^p_{\mathrm r}(G)$ implies the Powers property.
The result below does not prove that implication, but shows that simplicity
already rules out non-trivial amenable normal subgroups.

\section*{Source Question}

Phillips proves that $C^*_{\mathrm r}(G)$ simple implies that
$F^p_{\mathrm r}(G)$ and $B^{p,*}_{\mathrm r}(G)$ have the Powers property and
hence are simple. In the introduction, around source lines 408--414, he writes
that the reverse implications are not addressed and that it is not known
whether simplicity of $F^p_{\mathrm r}(G)$ implies the Powers property.

\section*{Idea of the Proof}

Let $N$ be an amenable normal subgroup of $G$. The quotient representation on
$\ell^p(G/N)$ can be approximated by the regular representation on
$\ell^p(G)$: spread a vector on $G/N$ uniformly over a finite Folner set in
each coset. For a fixed finitely supported element of $\C[G]$, only finitely
many translations inside $N$ are relevant, and amenability makes the Folner
set almost invariant under all of them. This proves that the quotient map on
group rings is contractive for the reduced $L^p$ norms.

\section*{Theorem}

\begin{theorem}\label{thm:quotient}
Let $G$ be a discrete group, let $N$ be an amenable normal subgroup, and let
$q:G\to G/N$ be the quotient map. Let $1\leq p<\infty$. Then the induced
homomorphism
\[
  \pi:\C[G]\to \C[G/N], \qquad \pi(g)=q(g),
\]
is contractive for the reduced $L^p$ operator norms. Hence it extends to a
contractive homomorphism with dense range
\[
  \Pi_N:F^p_{\mathrm r}(G)\longrightarrow F^p_{\mathrm r}(G/N).
\]
\end{theorem}

\begin{proof}
Write $\lambda_{G,p}$ and $\lambda_{G/N,p}$ for the left regular
representations on $\ell^p(G)$ and $\ell^p(G/N)$. It suffices to prove that
\[
  \|\lambda_{G/N,p}(\pi(a))\| \leq \|\lambda_{G,p}(a)\|
  \qquad (a\in \C[G]).
\]

Fix $a=\sum_{g\in E} a_g g\in\C[G]$ with $E$ finite. Choose a section
$\sigma:G/N\to G$ of $q$. For $g\in E$ and $t\in G/N$, define
\[
  \alpha(g,t)=\sigma(q(g)t)^{-1}g\sigma(t)\in N .
\]
First let $\xi\in \ell^p(G/N)$ have finite support $L$. For a finite non-empty
set $F\subset N$, define an isometry $V_F:\ell^p(G/N)\to\ell^p(G)$ by
\[
  (V_F\xi)(\sigma(t)n)= |F|^{-1/p}\xi(t)\mathbf 1_F(n)
  \qquad (t\in G/N,\ n\in N).
\]
For a basis vector $\delta_t$ and $g\in E$, the vectors
$\lambda_{G,p}(g)V_F\delta_t$ and $V_F\lambda_{G/N,p}(q(g))\delta_t$ are both
supported in the coset $\sigma(q(g)t)N$, and their difference has norm
\[
  \left(\frac{|\alpha(g,t)F\triangle F|}{|F|}\right)^{1/p}.
\]
Since $N$ is amenable and the set of all $\alpha(g,t)$ with $g\in E$ and
$t\in L$ is finite, there is a net of finite non-empty sets $F$ such that all
these quantities tend to zero. Therefore
\[
  \|\lambda_{G,p}(a)V_F\xi -
       V_F\lambda_{G/N,p}(\pi(a))\xi\| \longrightarrow 0 .
\]
Using that $V_F$ is an isometry, we get
\[
\begin{aligned}
  \|\lambda_{G/N,p}(\pi(a))\xi\|
  &= \lim_F \|V_F\lambda_{G/N,p}(\pi(a))\xi\| \\
  &\leq \limsup_F \|\lambda_{G,p}(a)V_F\xi\|
   \leq \|\lambda_{G,p}(a)\|\,\|\xi\|.
\end{aligned}
\]
Finite-support vectors are dense in $\ell^p(G/N)$, so the desired operator
norm inequality follows. The extension to completions is then immediate, and
its range contains $\C[G/N]$, hence is dense in $F^p_{\mathrm r}(G/N)$.
\end{proof}

\begin{corollary}\label{cor:obstruction}
Let $1<p<\infty$. If $G$ has a non-trivial amenable normal subgroup, then
$F^p_{\mathrm r}(G)$ is not simple. The same is true for
$B^{p,*}_{\mathrm r}(G)$.
\end{corollary}

\begin{proof}
Let $1\neq n\in N$. The map $\Pi_N$ sends the canonical generator $w_n$ to
$1$, so $w_n-1\in\ker \Pi_N$. This element is nonzero in
$F^p_{\mathrm r}(G)$: the canonical trace $\tau_p$ satisfies
$\tau_p(w_n)=0$ and $\tau_p(1)=1$. The kernel is proper because
$\Pi_N(1)=1$. Thus $F^p_{\mathrm r}(G)$ has a nonzero proper closed ideal.

For $B^{p,*}_{\mathrm r}(G)$, the norm on $\C[G]$ is the maximum of the
reduced $p$-norm of $a$ and the reduced $p$-norm of $a^*$. Since
$\pi(a^*)=\pi(a)^*$, Theorem~\ref{thm:quotient} gives the same contractive
dense-range homomorphism, and the same kernel argument applies.
\end{proof}

\begin{corollary}
If $1<p<\infty$ and either $F^p_{\mathrm r}(G)$ or
$B^{p,*}_{\mathrm r}(G)$ is simple, then the amenable radical of $G$ is
trivial.
\end{corollary}

\section*{Limitations}

This is only a necessary condition. It does not prove that simplicity of
$F^p_{\mathrm r}(G)$ implies the Powers property, nor that it implies
$C^*_{\mathrm r}(G)$ is simple. Possible remaining counterexamples must have
trivial amenable radical; the dynamical obstructions to $C^*$-simplicity are
not addressed here.

\section*{Relation to Gardella--Thiel}

Gardella--Thiel, arXiv:1408.6137, proved a quotient theorem for
$F^p_\lambda(G)\to F^p_\lambda(G/N)$ when $N$ is amenable and $G/N$ is finite,
and noted that the finite quotient assumption was likely unnecessary. The
Folner-spreading proof above gives the contractive dense-range map without
that finite quotient assumption, which is enough for the simplicity obstruction
used here.

\section*{References}

\begin{itemize}
\item N. Christopher Phillips, \emph{Simplicity of reduced group Banach
algebras}, arXiv:1909.11278v1.
\item E. Gardella and H. Thiel, \emph{Functoriality of group algebras acting
on $L^p$-spaces}, arXiv:1408.6137.
\end{itemize}

\end{document}
